\protect \section *{\nameref *{sec:plasma-scales}}
\begin{Solution}{12.{1}}
		іоносфера: $\lambda_D\approx0{,}7$~см, $a\approx2\times10^{-2}$~см, $N_D\approx2\times10^{5}$;
корона: $0{,}7$~см, $2\times10^{-3}$~см, $2\times10^{8}$; токамак: $7\times10^{-3}$~см, $2\times10^{-5}$~см, $2\times10^{8}$;
ядро Сонця: $\lambda_D\approx2\times10^{-9}$~см${}\sim a$, $N_D\approx5$~--- колективний опис найменш обґрунтований.
	
\end{Solution}
\begin{Solution}{12.{2}}
		$n^{-1/3}\approx2\times10^{-5}$~см${}<r_{Le}\approx4{,}8\times10^{-3}$~см${}<\lambda_{De}\approx7{,}4\times10^{-3}$~см${}<r_{Li}\approx0{,}29$~см;
$\lambda_{De}$ трохи більша за~$r_{Le}$.
	
\end{Solution}
\begin{Solution}{12.{3}}
		$\nu_{ei}\approx1{,}4\times10^{6}~\mathrm{с}^{-1}$, $\nu_{ei}/\omega_{pe}\approx2{,}5\times10^{-5}$,
$\lambda_{\mathrm{пр}}\approx1$~м; $\lambda_{\mathrm{пр}}/\lambda_{De}\approx4\times10^{4}$.
	
\end{Solution}
\begin{Solution}{12.{4}}
		$\lambda_D\approx1{,}3\times10^{-2}$~см, $\omega_{pe}\approx5{,}6\times10^{9}~\mathrm{с}^{-1}$ ($f\approx0{,}9$~ГГц);
обидва масштаби набагато менші за~сантиметрові розміри й~мікро-/мілісекундні часи розряду, тому плазма квазінейтральна.
	
\end{Solution}
\begin{Solution}{12.{5}}
		(а)~$\sim10^{5}$~с, (б)~$\sim10^{-2}$~с.
	
\end{Solution}
\begin{Solution}{12.{6}}
		$\Gamma\sim1$ при $n\sim3\times10^{29}~\mathrm{см}^{-3}$, $T\sim E_F$ при $n\sim10^{26}~\mathrm{см}^{-3}$.
	
\end{Solution}
\begin{Solution}{12.{7}}
		$Kn\approx10^{4}$, $r_{Li}/L\approx3\times10^{-3}$; обмеження показує $Kn$: поздовжній перенос
беззіткнювальний, тиск анізотропний~--- потрібні модель CGL, дрейфово- чи гірокінетичний опис поздовжнього переносу.
	
\end{Solution}
